\documentclass[10pt,a4paper]{amsart}
\usepackage[margin=19mm]{geometry}
\usepackage{amsmath,amssymb,mathtools}
\usepackage[scr=rsfs,cal=euler]{mathalfa}
\usepackage{xfrac}
\usepackage[unicode,hidelinks,bookmarks=false]{hyperref}
\newcommand{\ie}{i.e.\ }
\newcommand{\cf}{cf.\ }
\newcommand{\resp}{respectively\ }
\newcommand{\Subsection}{\subsection}
\providecommand{\nobreakdash}{\nobreak\hbox{-}}
\setlength{\emergencystretch}{2em}
\multlinegap0pt
\usepackage{bm}
\usepackage{bbm}
\usepackage{dsfont}
\usepackage{xspace}
\usepackage{cancel}
\usepackage{mathtools}
\usepackage{amsthm}
\makeatletter
\@ifundefined{thm}{\newtheorem{thm}{Thm}}{}
\@ifundefined{prop}{\newtheorem{prop}[thm]{Proposition}}{}
\makeatother
\newcommand{\bP}{\ensuremath{{\mathbb{P}}}}
\newcommand{\fo}{\ensuremath{{\mathfrak{o}}}}
\newcommand{\fp}{\ensuremath{{\mathfrak{p}}}}
\title[Finite index theorems for iterated Galois groups of preperio]{Finite index theorems for iterated Galois groups of preperiodic points for unicritical polynomials}
\author{Minsik Han, Thomas J. Tucker}
\date{Source: arXiv:2508.00266v2 (CC BY 4.0)}
\begin{document}
\maketitle

\noindent\emph{Source and licence.} Finite index theorems for iterated Galois groups of preperiodic points for unicritical polynomials. Version arXiv:2508.00266v2 (CC BY 4.0), distributed under the Creative Commons Attribution 4.0 International licence, as recorded in the arXiv licence metadata for this exact version and shown on the abstract page https://arxiv.org/abs/2508.00266v2. Full rights notice: licenses/arxiv-2508.00266-CC-BY-4.0.txt.\par\smallskip

\noindent\textit{Editorial scope.} Counted unit: the Prop whose source label is \texttt{main-d} in the article (source0.tex lines 480--494, source label \texttt{main-d}), delivered together with the article's own complete original proof of it (source0.tex lines 496--564, one \texttt{proof} environment of 69 lines). No mathematics is written, repaired or re-derived: statement and proof are the article's own text line for line. Required context retained uncounted: none. Every definition, display and label of those lines is retained unchanged. Omitted from this package and not redistributed: 1--479, 495--495, 565--1257. Nothing inside a retained excerpt is omitted or altered: every retained body is delivered exactly as the article's lines are written, so no omission receipt of any kind accompanies this package. Citations: the retained excerpts carry no citation command at all, so no bibliography entry of the article is transcribed and no reference is redistributed. Reference adaptation: none was needed and none was made; every reference inside the delivered unit resolves inside it, so no unresolved reference or citation is produced. Printed slips kept literal: no printed word, symbol, hypothesis or number of the counted statement or of its proof was altered. Layout and class: the article's own class is \texttt{amsart} with packages including pdfsync, amssymb, amsthm, amsmath, enumerate, hyperref, color; the wrapper is the lane's pinned \texttt{amsart} editorial wrapper at 10pt on A4, which restates the theorem-like environments the retained text needs and adds the article's own macro definitions that the retained text needs (\texttt{\textbackslash bP}, \texttt{\textbackslash fo}, \texttt{\textbackslash fp}), with extra wrapper packages (bm, bbm, dsfont, xspace, cancel, mathtools). The delivered numbering is the unit's own. Assets: no retained span contains a figure environment, a graphics-inclusion command, a PDF-page inclusion, a table or a TikZ picture; no image file is copied into this package, the package declares no asset dependency and no image file is redistributed with it. Licence: version arXiv:2508.00266v2 of this article is distributed under the Creative Commons Attribution 4.0 International licence, as recorded in the arXiv licence metadata for this exact version and shown on the abstract page https://arxiv.org/abs/2508.00266v2. A full rights notice is delivered at licenses/arxiv-2508.00266-CC-BY-4.0.txt. Characters: every retained excerpt is pure ASCII, so the delivered engine, XeLaTeX with its default Latin Modern text font, reports no missing character.
\begin{prop}\label{main-d}
  Let $d > 1$ be an integer, let $K$ be a number field, let
  $f$ be a polynomial of degree greater than 1, and let
  $S$ be a finite set of primes of $K$.  Let
  $\alpha \in \bP^1(K)$ be a point that is not preperiodic for
  $f$.  Let $\beta \in K$ be a point that is strictly preperiodic for
  $f$ but not post-critical for $f$.  Then for all
  sufficiently large $n$, there is a prime $\fp \notin S$ such that
  \begin{enumerate}
    \item[(1)] $v_\fp(f^n(\alpha) - \beta)$ is positive and not
      divisible by $d$; and
      \item[(2)] $v_\fp(f^m(\alpha) - \beta) = 0$ for all $m < n$.
  \end{enumerate}
  
\end{prop}

\begin{proof}
  Since $\beta$ is strictly preperiodic, the set $\{ f(\beta), \dots,
  f^k(\beta), \dots \}$ is finite and does not include $\beta$ as an
  element.  Thus, there are at most finitely many primes $\fp$ of $K$
  such that $f^k(\beta) \equiv \beta \pmod {\fp}$ for some $k >
  0$, so after expanding $S$ to a larger finite set, we may assume
  that it contains all such primes $\fp$.  Likewise, also possibly
  after expanding $S$, we may assume that the ring $\fo_S$ of
  $S$-integers in $K$ (that is elements of $K$ that are integers at
  every prime outside of $S$) is a principal ideal domain.  We may
  also assume that $S$ contains all primes of bad reduction for $f$
  and all primes $\fp$ such that $v_\fp(\alpha) < 0$ or $v_\fp(\beta)
  < 0$.  

  Now, if $v_\fp(f^n(\alpha) - \beta)$ is positive  for $\fp
  \notin S$, then we cannot have $v_\fp(f^m(\alpha) - \beta) >
  0$ for any $m < n$ since if we did we would have $f^{n-m}(\beta)
  \equiv \beta \pmod{\fp}$, which is impossible since $\fp \notin
  S$. Similarly, we cannot have $v_\fp(f^m(\alpha) - \beta) <  0$ for
  any $m$ since $v_\fp(f^m(\alpha))$ and $v_\fp(\beta)$ are both always
  non-negative since $f$ has good reduction at $\fp$ and $\alpha$ and
  $\beta$ are both integers at $\fp$.  Thus, condition (2) above will
  be met whenever (1) is, so it suffices to show that there is an
  $n_0$ such that for all $n \geq n_0$, there is an $\fp \notin S$
  such that $v_\fp(f^n(\alpha) - \beta)$ is positive and not
  divisible by $d$.

  The ring of $\fo_S^*$ of $S$-units is a finitely generated group, so
  $\fo_S^* / (\fo_S^*)^d$ is a finite group.  Let
  $\gamma_1(\fo_S^*)^d, \dots, \gamma_N(\fo_S^*)^d$ be
  the set of cosets of $(\fo_S^*)^d$ in $\fo_S^*$.  Since
  $\beta$ is not post-critical, $f^3(x) - \beta$ has at least 8 and
  thus more than 4 roots.  It follows from Riemann-Hurwitz that for
  each $i$, the curve $C_i$ given by $y^d = \gamma_i (f^3(x) - \beta)$
  has genus greater than 1. By Faltings' theorem, this means that each
  $C_i$ has finitely many rational points.  Thus, since $\alpha$ is
  not preperiodic there is an $n_0$ such that for all $n \geq n_0$
  there is no $y \in K$ such that
  $y^d = \gamma_i (f^3(f^{n-3}(\alpha)) - \beta)$.

  Now, let $n \geq n_0$ and suppose that for every prime $\fp \in S$,
  we have that $v_\fp(f^n(\alpha) - \beta)$ is either 0 or divisible
  by $d$.  Then the $\fo_S$ ideal generated by $f^n(\alpha) - \beta$
  is the $d$-th power of the $\fo_S$-ideal $I$.  Let $z$ be a generator
  for $I$ as an $\fo_S$ ideal (such a $z$ exists since $\fo_S$ is a
  principal ideal domain).  Then we have
  $z^d = u (f^n(\alpha) - \beta)$ for a unit $u \in \fo_s^*$.  We may
  write $u = \gamma_i w^d$ for one of our coset representatives
  $\gamma_i$ and some unit $w \in \fo_S^*$.  Let $y = z/w$.  Then we
  have $y^d = \gamma_i (f^3(f^{n-3}(\alpha))) - \beta)$ with $y \in K$,
  a contradiction.

 Thus, for every $n \geq n_0$, there is an $\fp \notin S$ such that
 $v_\fp(f^n(\alpha) - \beta)$ is positive and not divisible by $d$.  


  

  

      
    % First, for any $f\in \textup{PGL}_2(\Kbar)$ such that $\infty$ is a fixed point of $\psi = f^{-1}\varphi f$, we have \[v_\fp (\varphi^n(\alpha)-\beta) = v_\fp (f ( \psi^n (f^{-1}(\alpha))) - f(f^{-1}(\beta))) = v_\fp (\psi^n (f^{-1}(\alpha))-f^{-1}(\beta))\] for any $n$. (??) Therefore, we may assume that $\infty$ is a fixed point of $\varphi$.
    
    % Next, we extend $\cS$ to contain every $\fp$ such that \[\varphi^k(\beta) \equiv \beta \pmod \fp\] for any $k>0$. Since $0$ is strictly preperiodic, there are finitely many distinct values of $\varphi^k(\beta)$ but any of those is not equal to $\beta$. Therefore, only finitely many primes are added to $\cS$. We also add to $\cS$ every $\fp$ such that $v_\fp(\beta)<0$, so that $v_\fp(\beta)\ge 0$ for all $\fp \notin \cS$.

    % Now \cite[Proposition 12]{BIJJ} implies that there is $\fp \notin \cS$ satisfying the condition (1).

    % If $v_\fp(\varphi^m(\alpha)-\beta)>0$ for some $m<n$, then \[\varphi^m(\alpha) \equiv \beta \equiv \varphi^n(\alpha) \pmod{\fp},\] so \[\varphi^{n-m} (\beta) \equiv \beta \pmod \fp.\] This is impossible since we included every such $\fp$ in $\cS$. On the other hand, if $v_\fp(\varphi^m(\alpha)-\beta)<0$ for some $m<n$, then $v_\fp(\varphi^n(\alpha)-\beta)<0$ as well, since we assumed that $\infty$ is a fixed point of $\varphi$ and $v_\fp(\beta)\ge 0$. Therefore, the $\fp$ also satisfies the condition (2).
\end{proof}

\end{document}
